\section{Formal semantics}
\label{sec:semantics}

\S\ref{sec:design} presents Yuho informally; this section pins the
operational semantics for the fragment of the language the
empirical claims rely on. We give small-step rules for element
evaluation, prioritised-exception precedence in the
Catala~\cite{catala2021} default-logic style, the cross-section
composition predicates \texttt{apply\_scope} and
\texttt{is\_infringed}, and the penalty algebra. The body of this
section is the two-page summary reviewers can follow without
flipping to the appendix; the complete inference rules
(including the elided cases for ``mens rea'' typed-intent
variants, temporal-ordering constraints, and the burden-of-proof
qualifier) appear in \S\ref{sec:limitations}. We
deliberately omit any mechanisation here --- the soundness
theorem of \S\ref{sec:soundness} rests on the pen-and-paper rules
below; a Coq mechanisation is left as future work and noted in
\S\ref{sec:limitations}.

\subsection{Preliminaries}
\label{subsec:semantics-prelim}

\paragraph{Syntactic domains.} A Yuho \emph{module} $M$ is a
multiset of \emph{statutes}; a statute $S = \langle n, D,
\mathit{Es}, P, \mathit{Xs}, \mathit{Subs} \rangle$ where $n$ is
its section number (a string like \texttt{300} or
\texttt{376AA}), $D$ is a set of named definitions, $\mathit{Es}$
is a tree of element nodes and groups, $P$ is an optional penalty
algebra term, $\mathit{Xs}$ is a list of exception nodes carrying
priority and \texttt{defeats} edges, and $\mathit{Subs}$ is a
sequence of subsection sub-statutes.

An \emph{element} $e$ is a triple $\langle \kappa, x, \rho
\rangle$ where $\kappa \in \{\textsf{actus\_reus},
\textsf{mens\_rea}, \textsf{circumstance}\}$ tags the doctrinal
kind, $x$ is the element's identifier, and $\rho$ is its
free-text description (semantically opaque to the rules below).
An \emph{element group} $g = \langle c, \vec{m} \rangle$ has a
combinator $c \in \{\textsf{all\_of}, \textsf{any\_of}\}$ and an
ordered sequence of member nodes (each a leaf $e$ or a nested
group). \emph{Exception} $\xi = \langle \ell, \gamma, \delta,
\pi \rangle$ has a label $\ell$, a guard expression $\gamma$
(a boolean predicate over facts), an optional set $\delta$ of
sibling-exception labels it defeats, and an optional priority
$\pi \in \mathbb{N}$ (lower is higher-priority).

\paragraph{Semantic domains.} A \emph{fact pattern} $F$ is a
finite map from element identifiers to truth values, augmented
with the typed constants the encoded library uses
(\textsf{Intent} enums, money quantities, durations). The
\emph{statute environment} $\Sigma$ is a finite map from section
numbers to statute terms; given $M$, $\Sigma_M = \{\, n \mapsto S
\mid S \in M, S = \langle n, \ldots \rangle \,\}$. Evaluation
takes place under a context $\Gamma = \langle F, \Sigma \rangle$.

\paragraph{Judgements.} We define five judgements, in increasing
order of compositional depth:

\begin{itemize}
  \item $\Gamma \vdash e \Downarrow b$ — element $e$ evaluates to
    boolean $b$ under context $\Gamma$.
  \item $\Gamma \vdash g \Downarrow b$ — element group $g$
    evaluates to $b$.
  \item $\Gamma \vdash \xi \Downarrow_{\!\textsc{f}} b$ —
    exception $\xi$ \emph{fires} as $b$.
  \item $\Gamma \vdash S \Downarrow_{\!\textsc{c}} b$ — section
    $S$ produces a conviction outcome $b$.
  \item $\Gamma \vdash p \Downarrow_{\!\textsc{p}} \mathcal{R}$ —
    penalty term $p$ reduces to a sentence range $\mathcal{R}$.
\end{itemize}

We elide explicit threading of $\Gamma$ when it does not change
across rule premises. All judgements are \emph{total}: the rules
below cover every well-formed input, and the encoded library's
type-checker (\S\ref{sec:design}) rejects any module that would
land outside this coverage.

\subsection{Element evaluation}
\label{subsec:semantics-elements}

The element-kind tag $\kappa$ is doctrinally meaningful but
semantically inert at the structural layer: under
\textsc{Eval-Elem} below, all three kinds reduce by the same fact
lookup. The downstream verifier
(\S\ref{sec:implementation} Z3 backend) refines this with intent
typing for $\textsf{mens\_rea}$ elements; that refinement is
captured in \S\ref{sec:limitations}.

\begin{equation*}
\small
\inferrule[Eval-Elem]
  {F(x) = b}
  {\langle F, \Sigma \rangle \vdash \langle \kappa, x, \rho
   \rangle \Downarrow b}
\end{equation*}
\begin{equation*}
\small
\inferrule[Eval-Elem-Missing]
  {x \notin \mathrm{dom}(F)}
  {\langle F, \Sigma \rangle \vdash \langle \kappa, x, \rho
   \rangle \Downarrow \bot}
\end{equation*}

\textsc{Eval-Elem-Missing} is the key rule a reviewer should
notice: if a fact pattern omits an element identifier the section
asks about, the element evaluates to $\bot$. This is the
fail-closed reading the encoder relies on: the prosecution must
supply each element, not the encoder.

Element groups distribute over their member judgements:

\begin{equation*}
\small
\inferrule[Eval-AllOf]
  {\Gamma \vdash m_1 \Downarrow b_1 \quad \cdots \quad
   \Gamma \vdash m_k \Downarrow b_k}
  {\Gamma \vdash \langle \textsf{all\_of}, \vec{m} \rangle
   \Downarrow b_1 \wedge \cdots \wedge b_k}
\end{equation*}

\begin{equation*}
\small
\inferrule[Eval-AnyOf]
  {\Gamma \vdash m_1 \Downarrow b_1 \quad \cdots \quad
   \Gamma \vdash m_k \Downarrow b_k}
  {\Gamma \vdash \langle \textsf{any\_of}, \vec{m} \rangle
   \Downarrow b_1 \vee \cdots \vee b_k}
\end{equation*}

The top-level element tree of a statute is implicitly
$\textsf{all\_of}$. We write $\Gamma \vdash \mathit{Es}
\Downarrow_{\!\textsc{el}} b$ for the resulting truth value of
the section's element conjunction.

\subsection{Exception precedence}
\label{subsec:semantics-exceptions}

The Catala~\cite{catala2021} default-logic pattern: each
exception carries an explicit guard and is ordered by an explicit
priority DAG. Yuho extends this slightly: in addition to a
numeric \texttt{priority}, an exception may declare
\texttt{defeats $\ell'$} naming a sibling it overrides directly
(\S\ref{subsec:exceptions}). Both forms collapse into a single
``$\xi'$ defeats $\xi$'' relation, and we require its transitive
closure to be acyclic --- the linter
(\S\ref{sec:implementation}) enforces this and rejects modules
that would otherwise be ambiguous.

An exception's guard $\gamma$ is itself a boolean expression
over facts; its truth value $\Gamma \vdash \gamma \Downarrow b$
is given by standard expression evaluation
(\S\ref{sec:limitations}). The exception ``fires'' iff
its guard is true \emph{and} no defeating sibling also fires:

\begin{equation*}
\small
\inferrule[Fires]
  {\Gamma \vdash \gamma \Downarrow \top \\\\
   \forall \xi' \in \mathrm{defeats}^{-1}(\xi) :
   \Gamma \vdash \xi' \Downarrow_{\!\textsc{f}} \bot}
  {\Gamma \vdash \xi \Downarrow_{\!\textsc{f}} \top}
\end{equation*}

\begin{equation*}
\small
\inferrule[NotFires-Guard]
  {\Gamma \vdash \gamma \Downarrow \bot}
  {\Gamma \vdash \xi \Downarrow_{\!\textsc{f}} \bot}
\end{equation*}
\begin{equation*}
\small
\inferrule[NotFires-Defeated]
  {\Gamma \vdash \gamma \Downarrow \top \\
   \exists \xi' \in \mathrm{defeats}^{-1}(\xi) :
   \Gamma \vdash \xi' \Downarrow_{\!\textsc{f}} \top}
  {\Gamma \vdash \xi \Downarrow_{\!\textsc{f}} \bot}
\end{equation*}

The mutually-recursive shape of \textsc{Fires} on the priority
DAG is well-founded by construction (the linter rejects cycles).
A section's \emph{conviction} judgement closes the loop: elements
all hold \emph{and} no exception fires.

\begin{equation*}
\small
\inferrule[Conviction]
  {\Gamma \vdash \mathit{Es} \Downarrow_{\!\textsc{el}} \top \\\\
   \forall \xi \in \mathit{Xs} :
   \Gamma \vdash \xi \Downarrow_{\!\textsc{f}} \bot}
  {\Gamma \vdash S \Downarrow_{\!\textsc{c}} \top}
\end{equation*}

\begin{equation*}
\small
\inferrule[NoConviction-Elements]
  {\Gamma \vdash \mathit{Es} \Downarrow_{\!\textsc{el}} \bot}
  {\Gamma \vdash S \Downarrow_{\!\textsc{c}} \bot}
\end{equation*}
\begin{equation*}
\small
\inferrule[NoConviction-Excused]
  {\Gamma \vdash \mathit{Es} \Downarrow_{\!\textsc{el}} \top \\
   \exists \xi \in \mathit{Xs} :
   \Gamma \vdash \xi \Downarrow_{\!\textsc{f}} \top}
  {\Gamma \vdash S \Downarrow_{\!\textsc{c}} \bot}
\end{equation*}

The split between \textsc{NoConviction-Elements} and
\textsc{NoConviction-Excused} is the doctrinal hinge: it
distinguishes ``elements not made out'' from ``elements met but
defendant excused''. The Z3 backend exploits this distinction by
introducing two named booleans, \texttt{<sX>\_elements\_satisfied}
and \texttt{<sX>\_conviction}, whose biconditional
\texttt{conviction $\Leftrightarrow$ elements\_satisfied $\land\,
\lnot$ any\_fires} is the formal correspondent of these three
rules taken together.

\subsection{Cross-section composition}
\label{subsec:semantics-cross-section}

Two predicates lift the conviction judgement to expression
position so that a parent statute can refer to a base offence
without restating its elements. The lam4-flavoured
\texttt{is\_infringed(sX)} answers ``does $sX$ convict under the
ambient facts'', and the Catala-flavoured
\texttt{apply\_scope(sX, F')} answers the same question against
an explicit (possibly different) fact pattern $F'$.

\begin{equation*}
\small
\inferrule[IsInfringed]
  {\Sigma(n) = S \\
   \langle F, \Sigma \rangle \vdash S \Downarrow_{\!\textsc{c}}
   b}
  {\langle F, \Sigma \rangle \vdash \texttt{is\_infringed}(n)
   \Downarrow b}
\end{equation*}

\begin{equation*}
\small
\inferrule[ApplyScope]
  {\Sigma(n) = S \\
   \langle F', \Sigma \rangle \vdash S \Downarrow_{\!\textsc{c}}
   b}
  {\langle F, \Sigma \rangle \vdash
   \texttt{apply\_scope}(n, F') \Downarrow b}
\end{equation*}

The two predicates differ only in the fact pattern they thread:
\textsc{IsInfringed} re-uses the ambient $F$ (lam4-style ``the
facts the parent caller already has''), \textsc{ApplyScope}
substitutes a parent-supplied $F'$ (Catala-style scope
composition with explicit input). The implementation
(\S\ref{sec:implementation} Z3 backend) realises both as the
same lazily-declared boolean atom
\texttt{<sX>\_conviction}, deduplicated by name --- so a
satisfying assignment of the parent's constraints carries the
inner section's conviction-truth as a first-class proposition.

\subsection{Penalty algebra}
\label{subsec:semantics-penalty}

Penalty terms reduce to a \emph{sentence range} $\mathcal{R}$, an
abstract value carrying optional imprisonment, fine, caning, and
death components, each as a pair $\langle \mathit{lo},
\mathit{hi} \rangle$ over the appropriate measurement domain
(days, dollars, strokes, or the discrete unit
$\{\textsf{death}\}$). The leaf rules wrap structural penalty
literals; the combinator rules compose them.

\begin{equation*}
\small
\inferrule[Pen-Imprisonment]
  { }
  {\Gamma \vdash \texttt{imprisonment}\,\mathit{lo}\,\texttt{..}\,
   \mathit{hi} \Downarrow_{\!\textsc{p}}
   \langle \mathit{lo}, \mathit{hi} \rangle_{\!\mathrm{imp}}}
\end{equation*}

The fine, caning, and death-penalty leaves are analogous and
omitted here. The two non-trivial combinators are
\texttt{cumulative} (every component of every child applies) and
\texttt{or\_both} (the sentencer picks one or both children); we
write $\sqcup$ and $\sqcap$ respectively for the corresponding
range operations:

\begin{equation*}
\small
\inferrule[Pen-Cumulative]
  {\Gamma \vdash p_1 \Downarrow_{\!\textsc{p}} \mathcal{R}_1
   \quad \cdots \quad
   \Gamma \vdash p_k \Downarrow_{\!\textsc{p}} \mathcal{R}_k}
  {\Gamma \vdash \texttt{cumulative}\{p_1, \ldots, p_k\}
   \Downarrow_{\!\textsc{p}}
   \mathcal{R}_1 \sqcup \cdots \sqcup \mathcal{R}_k}
\end{equation*}

\begin{equation*}
\small
\inferrule[Pen-OrBoth]
  {\Gamma \vdash p_1 \Downarrow_{\!\textsc{p}} \mathcal{R}_1
   \quad \cdots \quad
   \Gamma \vdash p_k \Downarrow_{\!\textsc{p}} \mathcal{R}_k}
  {\Gamma \vdash \texttt{or\_both}\{p_1, \ldots, p_k\}
   \Downarrow_{\!\textsc{p}}
   \mathcal{R}_1 \sqcap \cdots \sqcap \mathcal{R}_k}
\end{equation*}

The G8 (\texttt{fine := unlimited}) and G14
(\texttt{caning := unspecified}) refinements
(\S\ref{subsec:gaps}) appear in the leaves as the special
range-pair values $\langle 0, +\infty \rangle$ and
$\langle 0, ? \rangle$ respectively; the linter ensures these
sentinels do not propagate beyond the leaf via the combinators.

\subsection{What the rules do not cover}
\label{subsec:semantics-scope}

Three deliberate elisions, named here so the soundness theorem of
\S\ref{sec:soundness} can refer to them:

\textit{Subsection composition.} Subsections $\mathit{Subs}$ are
self-similar to the parent statute (each is a sub-statute with
its own elements / penalty / exceptions). The Z3 backend
(\S\ref{sec:implementation}) hoists subsection elements to the
top level when the parent has none, which is the encoding
choice the rules above implicitly assume; a fully compositional
treatment lives in \S\ref{sec:limitations}.

\textit{Definitional substitution.} The \texttt{definitions}
block introduces named string aliases that the
controlled-English transpiler expands. They are semantically
opaque at the conviction layer --- a definition does not
contribute to element evaluation --- and we treat them as
metadata.

\textit{Refinement-typed quantities.} Fines, durations, and
counts can carry refinement bounds (G8, G14). The rules above
treat the underlying numeric range; refinement-bound enforcement
is a type-check pass orthogonal to the small-step semantics, and
discharged before the rules above fire.

These elisions are honest about the gap between the rules in
this section and the encoded library's full surface. The
soundness claim that follows applies to the fragment defined
above; the fragment is precisely what the empirical claims of
\S\ref{sec:evaluation} (the Z3-driven scenario synthesis, the
case-law differential testing) exercise.
